\section{Alphabet barriers and incremental optimality checks}
\label{sec:lcs-bounds}

\subsection{The next bottleneck is unnecessary exact work}

The complete overlap implementation has a polynomial local bound, but some
compressed substring queries still exhaust millions of work units on simple
repeated words. For example, the longest common substring of
$(a^N b^2)^2$ and $(a^N c^2)^2$ has length $N$. A total multiplicity bound
on these doubled strings permits $2N$, so it does not prove that the first
match of length $N$ is optimal. Continuing with exact occurrence tables
can be much more expensive than discovering the match itself.

There is a stronger inexpensive obstruction: a common substring cannot
cross a signed letter that is absent from the other word. Here $b$ and $c$
are separators, and the longest possible common-alphabet run has length
$N$. The new implementation computes this bound on the grammar DAG and
uses it at both the root and individual grammar-rule pairs. It also
processes witnesses incrementally. Once a verified witness attains the
relevant upper bound, all remaining candidates for that scope are unnecessary.
This is exact optimality detection, not speculative termination.

\subsection{An exact summary of allowed-letter runs}

Let $C$ be the intersection of the two words' signed alphabets. For a
nonempty grammar word $W$, store three lengths: $p(W)$, its longest
prefix over $C$; $q(W)$, its longest suffix over $C$; and $r(W)$, its
longest contiguous substring over $C$. A terminal in $C$ has summary
$(1,1,1)$, and any other terminal has $(0,0,0)$. For a binary rule $W=UV$,
\[
\begin{aligned}
 p(W)&=p(U)+\mathbf 1_{p(U)=|U|}\,p(V),\\
 q(W)&=q(V)+\mathbf 1_{q(V)=|V|}\,q(U),\\
 r(W)&=\max\{r(U),r(V),q(U)+p(V)\}.
\end{aligned}
\]
The third formula includes the only possible run crossing the rule's cut.
All quantities are computed bottom up without expanding any word.

\begin{proposition}[Shared-alphabet upper bound]
For any two reachable grammar words $U,V$ in the query,
\[
 \operatorname{LCS}(U,V)\le\min(r(U),r(V)),
\]
where the summaries use the intersection $C$ of the original root alphabets.
The summaries take linear arithmetic work and storage in the reachable
grammar size, with the bit cost of stored lengths.
\end{proposition}
\begin{proof}
Induction on the grammar gives the stated prefix, suffix and maximum-run
meanings: a run either lies in one child or crosses the cut as a suffix
followed by a prefix. Every letter in a common substring of $U,V$ occurs
in both root words, so it lies in $C$. The substring is therefore contained
in an allowed-letter run on each side. The recurrence uses a constant
number of arithmetic operations per rule.
\end{proof}

The intersection is signed: a letter and its inverse are different string
symbols. Terminal node identifiers are unique for each signed letter in
one arena, so intersecting the reachable terminal sets computes $C$
exactly. If one root's entire alphabet is shared, every substring of that
root is over $C$; the implementation uses ordinary node lengths on that
side and skips its summary pass. Disjoint alphabets give an immediate
zero answer. Bounds are local to the query, since changing the other
word can change $C$; they are not permanent word-arena metadata.

This bound can be stronger than shared total counts. Counts constrain how
many matching letters exist, while the run bound also constrains their
contiguity. Neither changes the equality test or accepts a witness that
has not been checked.

\subsection{Using progress as soon as it is proved}

The complete search initially tries the existing exact prefix query. If
that does not finish the query, it builds root offsets and alphabet sets,
computes any needed run summaries, and reduces its global upper bound to
\[
 U=\min(\text{requested cap},r(X),r(Y)),
\]
with the requested cap omitted when absent. An already found prefix or
single-letter witness of length $U$ settles the query immediately.

Binary rules are examined in descending run-bound order. A rule with
bound no greater than the best known common substring length cannot
improve the answer, nor can any subsequent rule with a smaller bound.
For a selected pair $A,B$, its local upper bound is
$U_{A,B}=\min(r(A),r(B),U)$. The matcher now yields the zero-offset
extension immediately, followed by the critical progression extensions
as they are computed. After every witness it updates the best answer.
Reaching $U$ finishes the whole query; reaching $U_{A,B}$ stops that pair
and continues only where a larger answer is still possible.

Previously the matcher collected all extension witnesses for a pair
before considering any of them. Even a first witness attaining a known
upper bound could be followed by expensive occurrence tables and further
extensions. The new loop does not request those computations. A local
stop is never treated as global completion unless the global bound is
also attained. The original exact fallback remains available whenever
the inexpensive bounds are loose, including queries with identical signed
alphabets but different letter order.

\subsection{Restricted families with inexpensive exact completion}

The first family is $(a^N b^2)^2,(a^N c^2)^2$, with $N=2^h$. The common
alphabet is $\{a\}$, the maximum allowed run is $N$ on both sides, and
the initial prefix witness has length $N$. No cut-pair search is needed.

A second family separates discovery from stopping. Represent
$b a^N b$ and $c a^N c$ with root cuts in the middle of the long run:
\[
 (b a^{N/2})(a^{N/2}b),\qquad
 (c a^{N/2})(a^{N/2}c).
\]
There is no common prefix. The first aligned-cut extension matches
$N/2$ letters backwards and $N/2$ forwards, reaching the run bound $N$.
Thus the new loop finishes after one extension, without requesting the
suffix-prefix overlap sets. Computing bounds while retaining eager
witness collection does not obtain this same early stop.

On these particular binary-power parses, construction, offsets, summaries
and the required prefix/suffix comparisons use $O(h\log h)$ traversal
and arithmetic operations in the constant-time-dictionary model, on
$O(h)$-bit lengths. This includes grammar traversal and sorting. Allowing
also a conservative linear worst-case dictionary lookup overhead gives
a polynomial bit bound of $O(h^3\log(h+1))$. This statement concerns
these explicit families and parse shapes; the general LCS primitive
retains its polynomial fallback with larger costs.

There is also a direct relator-search example with a loose total-count
bound. Let $M=N/4$ and use
\[
 D=a^N b^2 a^M b^2,\qquad R=a^N c^2 a^M c^2.
\]
The shared count is $N+M$, but the maximum shared-alphabet run, including
cyclic seams, is $N$. The checked prefix attains it. For $N>M+4$, the
overlap gives positive raw gain $N-M-4$ and replaces the target by
\[
 b^{-2}a^{-M}b^{-2}c^2a^Mc^2.
\]
The total relator length drops from $2N+2M+8$ to $N+3M+12$.
The earlier whole-donor stage cannot match either relation because each
has a generator absent from the other. The new bound therefore accelerates
an actual partial-move query whose count cap alone is insufficient.
This is a supplied-presentation operation, not a completed knot verdict.

\subsection{Validation and controlled ablation}

Four new tests check run bounds against literal expansion at every reachable
node of 400 randomly parsed signed words; exponential uncapped queries
through $N=2^{500}$ with occurrence tables forbidden; an interior witness
with further overlaps forbidden; local-pair stopping against a literal
LCS oracle; a same-alphabet negative control; exhausted work; and the
large partial-move family above without expansion. The existing exhaustive
and randomized LCS, cyclic-overlap and witness tests are retained.

All 682 tests pass with Regina 7.4.1 in 182.222 seconds. The preceding
ten-test targeted run passed in 27.990 seconds; the integrated run also
checks the interior example with both root cuts inside the long run.

\begin{center}\small
\begin{tabular}{@{}lrrr@{}}
\toprule
Scope & Before & A/A & Both changes \\
\midrule
Small 14, summed medians & 280.857 & 280.086 & 282.904 \\
Gordian, full query & 1952.276 & 1939.691 & 1940.255 \\
\bottomrule
\end{tabular}
\end{center}
Median milliseconds; all whole-knot queries complete.


All 225 whole-knot queries retain identical certificate digests, and none
reaches the changed matcher. The small-fourteen sum of medians is 280.857,
280.086 and 282.904 milliseconds for baseline, A/A and both changes.
Gordian takes 1.952, 1.940 and 1.940 seconds respectively. These data check
preservation of existing paths; they do not demonstrate a knot-corpus speedup.

\begin{center}\footnotesize
\begin{tabular}{@{}llrrrrr@{}}
\toprule
LCS family & $h$ & Before & A/A & Bounds & Incr. & Both \\
\midrule
Repeated & 8 & 109.458 & 106.942 & 0.135 & 107.604 & 0.142 \\
Repeated & 32 & limit & limit & 0.245 & limit & 0.319 \\
Repeated & 500 & limit & limit & 3.078 & limit & 3.028 \\
Interior & 8 & 60.409 & 59.386 & 58.950 & 59.761 & 0.177 \\
Interior & 32 & limit & limit & limit & limit & 0.414 \\
Interior & 500 & limit & limit & limit & limit & 4.912 \\
Same alphabet & 8 & 1.124 & 1.116 & 1.111 & 1.127 & 1.117 \\
Same alphabet & 32 & 8.550 & 8.543 & 8.647 & 8.598 & 8.595 \\
Same alphabet & 500 & limit & limit & limit & limit & limit \\
\bottomrule
\end{tabular}
\end{center}
Median milliseconds including construction. Bounds retains eager witness collection; Incr. retains ordinary length bounds. Limits are censored queries.
\begin{center}\small
\begin{tabular}{@{}rrrrrr@{}}
\toprule
Partial move, $h$ & Before & A/A & Both & Nodes & Work \\
\midrule
8 & 287.444 & 286.253 & 0.409 & 35 & 527 \\
32 & limit & limit & 0.875 & 83 & 1439 \\
500 & limit & limit & 8.418 & 1019 & 19223 \\
\bottomrule
\end{tabular}
\end{center}
Median milliseconds including construction, failed whole-donor search, cyclic query and checked application. Nodes/work describe both changes.


The kernel audit contains 225 uncapped LCS measurements. At $h=8$, the
repeated family falls from 109.458 milliseconds to 0.142 milliseconds.
Bounds alone obtain this improvement; incremental witnesses alone do not.
At $h=500$, both changes complete in 3.028 milliseconds, with 509 nonempty
nodes and 5562 work units, while the original query exhausts its two-million-
work allowance.

The interior family separates the two effects. At $h=8$, archived search
takes 60.409 milliseconds, bounds alone 58.950, incremental witnesses alone
59.761, and both changes 0.177. At $h=500$, both changes finish in 4.912
milliseconds with 1012 nodes and 8562 work units; every isolated or archived
arm exhausts the cap. A bound does not save the later exact work unless
the checked witness is considered before that work is requested.

The same-alphabet control completes with identical work counts through
$h=32$: 16187 units and approximately 8.6 milliseconds in each arm.
At $h=500$, every arm still exhausts two million work units. Thus the new
bound leaves a measurable class of expensive queries unresolved at this
checkpoint. Section~\ref{sec:lcs-transitions} adds directed-pair constraints
that resolve this family, while retaining a stronger negative control.

Forty-five further measurements include the complete partial-move operation.
At $h=8$, construction plus search and application falls from 287.444 to
0.409 milliseconds. At $h=500$, the combined path completes in 8.418
milliseconds with 1019 nodes and 19223 work units; both archived arms hit
the cap. These are exact local operations on supplied relators, not
completed recognition of those groups. No speedup ratio is assigned to
an exhausted baseline.


The audit compares the archived matcher from \code{e678ef097} against
two isolated components and their combination. The bounds-only arm uses
run summaries and bound-based ordering but eagerly collects pair witnesses.
The incremental-only arm uses the new witness loop with the previous
ordinary-length bounds. An A/A arm repeats the archived class. These
classes run in the common current word arena, recognition and replay
harness, so this is a controlled component comparison.

One warmup is excluded before five measured rounds per case; arm order
is shuffled with seed 2910. Full knot timings include fresh planar-diagram
construction, recognition and independent replay; certificate digests are
checked outside the timer. Kernel timings include grammar construction
and uncapped exact LCS. A same-alphabet family,
$(ab)^Nac,(ab)^Nca$, has exact LCS length $2N$ and provides a control
where run summaries are skipped. The partial-move timings additionally
include the unsuccessful whole-donor query, cyclic search and checked
application. Resource exhaustion is censored, never a negative substring
answer or a completed-time denominator.

Whole recognition allows 18 seconds, with 15 seconds and twenty million
work units for the group stage. Kernels allow two million work units and
100000 nodes/table cells. Raw results are in
\path{results/lcs_bounds_20261008.json}; logs are in
\path{synthesis/data/lcs-bounds-*}. No new incoming report was present,
and the suggested \code{openai/math} head was unchanged at the previously
recorded \code{fd4aeeb2ee}. These are local refinements of the complete
compressed-overlap implementation. They do not settle global move count,
grammar growth or completeness, so the general subexponential recognition
bound remains unproved.
